\section{RIMap-RISC}

\begin{frame}{RIMap-RISC : idée clé}
  \vspace{2cm}
\begin{center}
\begin{minipage}{0.85\textwidth}
\begin{itemize}
	\item Outil de \textbf{prédiction transcriptome-wide} des interactions miRNA--transcrit.
	\item Développé par le laboratoire de \textbf{François Major} à l'IRIC.
	\item Modélise l'\textbf{architecture bipartite du complexe RISC}
	      (région seed + région supplémentaire + pont flexible).
\end{itemize}
\end{minipage}
\end{center}
\vfill
\scriptsize Chasles S, Gaillard-Duchassin Z, Quenneville J, Lemaire M, Gagnon E, Major F. RIMap-RISC: a transcriptome-wide database of structurally modeled human microRNA interactions. \textit{Genome Biology} 27, 112 (2026). \url{https://doi.org/10.1186/s13059-026-04008-y}
\end{frame}

\begin{frame}{Architecture bipartite de RISC}
\molochset{block=fill}
\begin{columns}[T, onlytextwidth]
	\column{0.55\textwidth}
	\begin{block}{Modèle RIMap-RISC}
		\begin{itemize}
			\item \textbf{Seed} (g2--g8) : au moins 4 paires Watson-Crick contiguës.
			\item \textbf{Pont} (g9--g12) : région flexible entre les deux domaines.
			\item \textbf{Supplémentaire} (g13--g17) : appariement additionnel optionnel.
		\end{itemize}
	\end{block}

	\column{0.40\textwidth}
	\vspace{0.5em}
	\begin{itemize}
		\item Supporte les \textbf{appariements non-canoniques} : mésappariements, wobble, bulges.
		\item Prédiction de structure : \textbf{MC-Flashfold}.
		\item Énergie libre $\Delta G$ et $K_d$ estimés à partir de données cinétiques (RNA Bind-n-Seq).
	\end{itemize}
\end{columns}
\vfill
\scriptsize Dallaire P, Major F. MC-Flashfold: a fast RNA secondary structure prediction engine using a probabilistic helix-based model. \textit{Methods Mol Biol} 1490:237--51 (2016). \url{https://doi.org/10.1007/978-1-4939-6433-8_15}
\end{frame}

\begin{frame}{RIMap-RISC : entrées et sorties}
\begin{columns}[T, onlytextwidth]
	\column{0.45\textwidth}
	\begin{block}{Entrée}
		Une paire \textbf{(transcrit, miRNA)} :
		\begin{itemize}
			\item Séquence du transcrit (région CDS / 3'UTR)
			\item Séquence du miRNA mature
		\end{itemize}
	\end{block}

	\column{0.50\textwidth}
	\begin{block}{Sortie : sites de liaison}
		Pour chaque site trouvé :
		\begin{itemize}
			\item Position sur le transcrit
			\item Type de seed (6mer, 7mer-m8, \ldots)
			\item $\Delta G$ (énergie libre)
			\item $K_d$ (constante de dissociation)
			\item Bridge, appariements
		\end{itemize}
	\end{block}
\end{columns}

\vspace{0.5em}
\centering
\small Une paire peut avoir \textbf{plusieurs sites de liaison}.
\end{frame}
